% !TEX program = lualatex
% あるいは pdfLaTeX を使う場合は CJKutf8 を有効にしてください。
\documentclass[12pt]{article}
\usepackage[margin=1in]{geometry}
\usepackage{amsmath,amssymb,amsthm,mathtools,bm}
\usepackage{physics}
\usepackage{tikz}
\usetikzlibrary{arrows.meta,calc,positioning}
\usepackage{xeCJK}
\usepackage{hyperref}
\usepackage[T1]{fontenc}
\usepackage{lmodern}
\usepackage{listings}
\lstset{basicstyle=\ttfamily\small,breaklines=true,columns=fullflexible,keepspaces=true}

\newcommand{\R}{\mathbb{R}}
\newcommand{\vct}[1]{\mathbf{#1}}
\newcommand{\cross}{\times}
\newcommand{\triple}[3]{\left[ #1,\,#2,\,#3 \right]}

\title{Space Vectors in \(\R^3\): Cross Product, Scalar Triple Product, and Jacobi Identity\\
空間ベクトル：外積・スカラー三重積・Jacobi 恒等式}
\author{Student (Your Name) \\
Co-author: CODEX (OpenAI)\\
\small TeX generated by CODEX (OpenAI); Lean file generated by CODEX (OpenAI)}
\date{September 15, 2025}

\begin{document}
\maketitle

\begin{abstract}
This handout summarizes definitions and core properties of the dot product, cross product, and the scalar triple product in \(\R^3\), and outlines a coordinate proof of the Jacobi identity for the cross product. Bilingual notes are provided (EN/JP), and the Lean formalization used for verification is excerpted.\\
本資料は \(\R^3\) における内積・外積・スカラー三重積の定義と基本性質を整理し、外積の Jacobi 恒等式の座標計算による証明概略を示します。検証に用いた Lean 形式化の抜粋も掲載します。
\end{abstract}

% =====================
\section*{1. Definitions / 定義}
\paragraph{English.}
For \(u=(u_1,u_2,u_3), v=(v_1,v_2,v_3)\in\R^3\):
\[
\langle u,v\rangle = u_1 v_1 + u_2 v_2 + u_3 v_3,\qquad
u\cross v = (u_2 v_3 - u_3 v_2,\; u_3 v_1 - u_1 v_3,\; u_1 v_2 - u_2 v_1).
\]
For \(w\in\R^3\), the scalar triple product is
\[
\triple{u}{v}{w} = \langle u,\, v\cross w\rangle = \det\begin{pmatrix}
u_1 & u_2 & u_3\\ v_1 & v_2 & v_3\\ w_1 & w_2 & w_3
\end{pmatrix}.
\]
\paragraph{日本語.}
\(u=(u_1,u_2,u_3), v=(v_1,v_2,v_3)\) に対し，
\[
\langle u,v\rangle = u_1 v_1 + u_2 v_2 + u_3 v_3,\qquad
u\cross v = (u_2 v_3 - u_3 v_2,\; u_3 v_1 - u_1 v_3,\; u_1 v_2 - u_2 v_1).
\]
また \(\triple{u}{v}{w} = \langle u, v\cross w\rangle\) は 3 本のベクトルが張る平行六面体の有符号体積（行列式）です。

% =====================
\section*{2. Worked Examples / 例題}
Let \(a=(1,2,3),\; b=(2,-1,1),\; c=(0,1,2),\; P=(1,1,1),\; n=(2,-1,3)\).
\begin{align*}
&\textbf{Q1: } a\cross b = (5,5,-5).\\
&\textbf{Q2: } u\cross v = - (v\cross u).\\
&\textbf{Q3: } \langle u,\, u\cross v\rangle = 0.\\
&\textbf{Q4: } \triple{a}{b}{c} = -5.\\
&\textbf{Q5: } (\langle (3,2,1) - P,\, n\rangle)^2 = 9.\quad (\text{ここで }\langle (3,2,1)-P, n\rangle = 3)
\end{align*}

\noindent\textit{JP 注記:} 問題文の平面点 \(P=(1,1,1)\) は式 \(2x-y+3z=6\) を満たしません（左辺=4）。距離の幾何と整合させるには，平面上の点へ修正するか \(|\langle Q,n\rangle-6|\) 型で扱うのが自然です。ここでは代数計算として \(\langle Q-P,n\rangle\) を評価しました。

% =====================
\section*{3. Core Properties / 基本性質}
\paragraph{English.}
(i) Anti-commutativity: \(u\cross v = -\,v\cross u\). (ii) Orthogonality: \(\langle u, u\cross v\rangle=0\) and \(\langle v, u\cross v\rangle=0\). (iii) Scalar triple product equals determinant, hence gives the oriented volume. (iv) \(|\triple{u}{v}{w}|\) equals the volume of the parallelepiped spanned by \(u,v,w\).

\paragraph{日本語.}
\(i)\) 反交換性：\(u\cross v = -\,v\cross u\)。\;\(ii)\) 直交性：\(\langle u, u\cross v\rangle=0\) および \(\langle v, u\cross v\rangle=0\)。\;\(iii)\) 三重積は行列式に等しく，有符号体積を与える。\;\(iv)\) 絶対値 \(|\triple{u}{v}{w}|\) は平行六面体の体積です。

% =====================
\section*{4. Jacobi Identity / Jacobi 恒等式}
\paragraph{Statement.}
\[
u\cross(v\cross w) + v\cross(w\cross u) + w\cross(u\cross v) = 0.
\]
\paragraph{Coordinate outline (EN).}
Using the Levi-Civita tensor \(\varepsilon_{ijk}\) and the identity
\(\varepsilon_{ijk}\varepsilon_{klm} = \delta_{il}\delta_{jm} - \delta_{im}\delta_{jl}\),
each component of the left-hand side becomes a cyclic sum of terms that cancel in pairs, yielding zero.

\paragraph{座標計算の概略（JP）.}
Levi-Civita 記号 \(\varepsilon_{ijk}\) と恒等式 \(\varepsilon_{ijk}\varepsilon_{klm}=\delta_{il}\delta_{jm}-\delta_{im}\delta_{jl}\) を用いると，各成分は巡回和で互いに打ち消し合い 0 になります。Lean では成分展開 → 多項式整理（`ring`）で機械検証しました。

% =====================
\section*{5. Lean Excerpts / Lean 抜粋}
Below are minimal Lean snippets used in the file \texttt{src/MyProjects/LinearAlgebra/A1/zenla1\_session02.lean}. We deliberately used function style `cross u v` and `scalar_triple u v w` to avoid notation conflicts.

\begin{lstlisting}
abbrev R3 := ℝ × ℝ × ℝ
def dot (u v : R3) : ℝ := u.1 * v.1 + u.2.1 * v.2.1 + u.2.2 * v.2.2
def cross (u v : R3) : R3 :=
  (u.2.1 * v.2.2 - u.2.2 * v.2.1,
   u.2.2 * v.1 - u.1 * v.2.2,
   u.1 * v.2.1 - u.2.1 * v.1)
def scalar_triple (u v w : R3) : ℝ := dot u (cross v w)

theorem Q1 : cross a b = (5, 5, -5) := by
  ext <;> norm_num [cross, a, b]

theorem Q2 (u v : R3) : cross u v = - (cross v u) := by
  ext <;> simp [cross] <;> ring

theorem Q3 (u v : R3) : ⟪u, cross u v⟫ = 0 := by
  simp [dot, cross]; ring

theorem Q4 : scalar_triple a b c = -5 := by
  norm_num [scalar_triple, dot, cross, a, b, c]

theorem Q5 : let Q : R3 := (3, 2, 1)
  (⟪Q - P, n⟫)^2 = 9 := by
  intro Q
  have hDot : ⟪Q - P, n⟫ = 3 := by
    norm_num [dot, Q, P, n]
  simpa only [hDot] using (by norm_num : (3:ℝ)^2 = 9)

theorem Challenge (u v w : R3) :
  cross u (cross v w) + cross v (cross w u) + cross w (cross u v) = 0 := by
  rcases u with ⟨ux, uy, uz⟩; rcases v with ⟨vx, vy, vz⟩; rcases w with ⟨wx, wy, wz⟩
  ext <;> simp [cross] <;> ring
\end{lstlisting}

% =====================
\section*{6. Figures / 図}
\subsection*{(a) Cross Product Direction / 外積の向き}
\begin{center}
\begin{tikzpicture}[scale=1.0,>=Latex]
  % axes
  \draw[->,thin,gray] (-0.2,0) -- (3.2,0) node[below] {$x$};
  \draw[->,thin,gray] (0,-0.2) -- (0,2.4) node[left] {$y$};
  % vectors u and v
  \draw[->,ultra thick,blue] (0,0) -- (2.2,0.8) node[below right] {$\,u$};
  \draw[->,ultra thick,orange] (0,0) -- (0.8,1.8) node[left] {$v\,$};
  % cross product (out of plane)
  \draw[->,ultra thick,green!60!black] (0,0) -- (0,0.0) coordinate (tmp);
  \draw[fill=green!60!black] (0,0) circle (1pt);
  \node[above right=2pt] at (0,0) {$u\cross v$};
  % wedge sector to indicate right-hand rule
  \draw[semithick,gray!70] (0:0.8) arc (0:65:0.8);
  \node at (0.55,0.35) {\small $\theta$};
\end{tikzpicture}
\end{center}

\subsection*{(b) Parallelepiped Volume / 平行六面体の体積}
\begin{center}
\begin{tikzpicture}[scale=1.0,>=Latex]
  % Oblique basis
  \coordinate (O) at (0,0);
  \coordinate (V) at (2.0,0.5);
  \coordinate (W) at (0.7,1.6);
  \coordinate (U) at (0.3,2.0);
  % base parallelogram spanned by v and w
  \fill[orange!20] (O) -- (V) -- ($(V)+(W)$) -- (W) -- cycle;
  % top face shifted by u
  \fill[orange!10] ($(O)+(U)$) -- ($(V)+(U)$) -- ($(V)+(W)+(U)$) -- ($(W)+(U)$) -- cycle;
  % edges
  \draw[->,ultra thick,blue] (O) -- (U) node[left] {$u$};
  \draw[->,ultra thick,black] (O) -- (V) node[below right] {$v$};
  \draw[->,ultra thick,black] (O) -- (W) node[left] {$w$};
  % vertical edges
  \draw[densely dashed] (V) -- ($(V)+(U)$);
  \draw[densely dashed] (W) -- ($(W)+(U)$);
  \draw[densely dashed] ($(V)+(W)$) -- ($(V)+(W)+(U)$);
  % label
  \node[above right] at ($(W)!0.5!(V)$) {$|\triple{u}{v}{w}|$};
\end{tikzpicture}
\end{center}

% =====================
\section*{7. How to Compile / コンパイル方法}
\begin{itemize}
  \item Recommended: \textbf{LuaLaTeX} or \textbf{XeLaTeX} (UTF-8, Japanese ready).
  \item With pdfLaTeX: enable \verb|CJKutf8| and wrap Japanese parts with the CJK environment (already done here).
\end{itemize}

\bigskip
\noindent\textbf{Acknowledgment.} Both the Lean file and this TeX were generated and curated with \textbf{CODEX (OpenAI)} using the Codex CLI workflow.\\
（謝辞）Lean ファイルおよび本 TeX は \textbf{CODEX (OpenAI)} により Codex CLI ワークフローで生成・整備されました。

\end{document}

